\documentclass[a4paper,12pt]{article}

% Font configuration for Times New Roman
\usepackage{mathptmx}
\usepackage[T1]{fontenc}

% Page geometry settings (25mm top/bottom, 20mm left/right)
\usepackage[top=25mm, bottom=25mm, left=20mm, right=20mm]{geometry}

% Remove page numbers, headers, and footers
\pagestyle{empty}

% Optional packages for tables and graphics
\usepackage{graphicx}
\usepackage{booktabs}

% Custom command for run-in headings
\newcommand{\heading}[1]{%
  \noindent\textbf{#1}%
}

\setlength{\parskip}{1em}
\setlength{\parindent}{0em}

\begin{document}

% Top Matter
\begin{center}
    {\fontsize{13.4pt}{12pt}\selectfont 
    \textbf{Abstract Template for Methods and Techniques in Phonetics Sciences (MaTiPS) 2026}\par}
    % \vspace{1em}
    {\fontsize{12pt}{10pt}\selectfont
    Author Full Name$^{1}$, Co-author Full Name$^{2}$\par
    \vspace{0.25em}
    $^{1}$Affiliation\\
    $^{2}$Affiliation\par
    \vspace{0.25em}
    Author 1 email address, Author 2 email address\par
    \vspace{-0.75em}
    }
\end{center}

\heading{General format.}
This is the template for MaTiPS 2026 conference abstracts. Abstracts may be up
to 2 A4-size pages long, including text, figures, tables, and references. The
top matter (title, authors, affiliations, email addresses) must be formatted as
in this template (the number of lines needed for the top matter will depend on
the length of the title and the number of authors). The beginning of the
abstract must be separated from the top matter by a blank line. All text should
be in Times New Roman font. The title is 14-point font, all other text is
12-point font. Margins are 25 mm (.98") top and bottom, and 20 mm (.79") left
and right. Footers and headers are not allowed. Do not use page numbers. 

All abstracts must be submitted in PDF format in compliance with the provided
template.The PDF file name should be structured as follows: 

\noindent\texttt{MaTiPS2026\_<presentation preference, i.e. Poster or Tutorial>\_<1st author's last name>\_<short title without spaces, e.g. AbstractTemplate>.pdf}

For example: \texttt{MaTiPS2026\_Poster\_Lulich\_AbstractTemplate.pdf}

\heading{Headings.}
Headings are in bold font with the first word capitalized and the rest of the
heading in lower case. Text follows without line break. No sub-headings.

\heading{Figures and tables.}
Figures and tables should be of high visual quality. No screenshots. Tables
should be created directly within the template.



\setlength{\leftmargin}{1.5em}%
\setlength{\itemindent}{-1.5em}%
\setlength{\parsep}{0pt}%
\setlength{\itemsep}{0.2em}%

\noindent\textbf{References}\\
References should start on the line immediately after the {\bf  References} header, and should be
formatted with hanging indentation. Use a consistent format for references and inline citations.

% Custom environment for hanging indentation references
% \newenvironment{reflist}{%
%   \begin{list}{}{%
%     \setlength{\leftmargin}{1.5em}%
%     \setlength{\itemindent}{-1.5em}%
%     \setlength{\parsep}{0pt}%
%     \setlength{\itemsep}{0.2em}%
%   }
% }{%
%   \end{list}
% }
% \begin{reflist}
% \item Author, A. A. (2026). Title of reference. \textit{Journal Name}, Volume, Page range.
% \item Use a consistent format for references and inline citations.
% \end{reflist}

\end{document}
